
\documentclass[%
	BCOR12mm, % Bindekorrektur 12mm (geeignet für Lochung oder einfache Klebebindung)
	cleardoublepage=empty, % Vakatseiten ohne Kopf- und Fußzeile
	tablecaptionabove, % Tabellenüberschriften, aber Bildunterschriften
	toc=bib, % Literaturverzeichnis im Inhaltsverzeichnis aufführen
	toc=listofnumbered, % Abbildungs- und Tabellenverzeichnis im Inhaltverzeichnis aufführen
	listof=leveldown, % Abbildungs- und Tabellenverzeichnis Überschriften als Section statt Chapter
	numbers=noendperiod % Kein abschließender Punkt bei Kapitel-Zählung
	]{scrbook}

\usepackage[utf8]{inputenc} % UTF-8 Zeichen sind möglich (z.B. ö, ü, ä, ...)
\usepackage[ngerman]{babel} % für deutsche Texte

\usepackage{amssymb,amsmath} % Matheumgebung, siehe "mathe_in_latex.pdf"
\usepackage{siunitx} % siehe "siunitx_Einheiten.pdf"
\usepackage{mhchem} % chemische Summenformeln und Gleichungen siehe "Summenformeln und Chemie allgemein mit LaTeX.pdf"
\usepackage{chemfig} % chemische Sturkturformeln darstellen siehe "Strukturformeln mit chemfig.pdf"
\usepackage{pgfplots} % Diagramme pltten siehe "pgfplots_Diagramme.pdf"
\usepackage{tikz} % zeichnen, Deckblattgestaltung, etc. siehe "tikzpgfmanual.pdf"
\usetikzlibrary{backgrounds}
\usepackage{hyperref}
\usepackage{color}

\usepackage{natbib} %Literaturverzeichnis

%% ----------------------------------
%% |  Style of appendix numbering   |
%% ----------------------------------
\renewcommand\appendix{\par
	\addchap{\appendixname}
	\setcounter{section}{0}% 
	\setcounter{subsection}{0}% 
	\setcounter{figure}{0}%
	\renewcommand\thesection{\Alph{section}}% 
	\renewcommand\thefigure{\Alph{section}.\arabic{figure}} 
	\renewcommand\thetable{\Alph{section}.\arabic{table}}}
%% --- End of appenix numbering style ---

%--------------------------------------------------------------------------------
\begin{document}
%--------------------------------------------------------------------------------
%	DECKBLATT
%--------------------------------------------------------------------------------
%--------------------------------------------------------------------------------
%	DECKBLATT
%--------------------------------------------------------------------------------
\thispagestyle{empty}

\begin{tikzpicture}[remember picture,overlay]
	
	\fill[color=kit-black!15] (current page.south east) rectangle (current page.north west);
	\fill[color=white] ($ (current page.north west) + (19mm,-5mm) $) -- ($ (current page.north east) - (10mm,5mm) $) arc (90:0:5 mm) -- ($ (current page.south east) + (-5mm,13mm) $) -- ($ (current page.south west) + (24mm,13mm) $) arc (270:180:5 mm) -- cycle;

	\node[anchor=north west] at ($ (current page.north west) + (26.5mm,-12.5mm) $){ \includegraphics[height=18mm]{../Logos/kitlogo_de_cmyk.pdf} };
	\node[anchor=north east] at ($ (current page.north east) - (12.5mm,12.5mm) $){ \includegraphics[height=18mm]{../Logos/LVTLogo.pdf} };
	
	\node[text width=0.9\textwidth] at ($ (current page.center) + (14mm,0mm) $){
		\titlehead{\Large \center Karlsruher Institut für Technologie (KIT)} % Universität, Institut, ...
		\subject{Abschlussarbeit} % Art der Arbeit (Bachelor-, Master-, Diplomarbeit, Dissertation)
		\title{Titel}
		\subtitle{Untertitel}
		\author{Autor}
		% \publishers{Aufgabensteller \\Betreuer}  % Aufgabensteller / Betreuer
		%\date{} %Datum wird automatisch auf \today gesetzt - anderes Datum in die {} schreiben
	\maketitle };
	
	\node[anchor=base west] at ($ (current page.south west) + (24mm,7mm) $) {Karlsruher Institut für Technologie (KIT)};
	\node[anchor=base east] at ($ (current page.south east) + (-10mm,7mm) $) {www.kit.edu};
	
\end{tikzpicture}
%--------------------------------------------------------------------------------
%--------------------------------------------------------------------------------
%	ERKLÄRUNG - INHALTSVERZEICHNIS - SYMBOLVERZEICHNIS
%--------------------------------------------------------------------------------
\frontmatter
%--------------------------------------------------------------------------------
% \chapter{Erklärung} % (fold)
% \label{cha:erklaerung}

% chapter erklärung (end)
%--------------------------------------------------------------------------------
\tableofcontents
%--------------------------------------------------------------------------------
\chapter{Symbolverzeichnis} % (fold)
\label{cha:symbolverzeichnis}

% chapter symbolverzeichnis (end)

%--------------------------------------------------------------------------------
%	INHALT
%--------------------------------------------------------------------------------
\mainmatter
%--------------------------------------------------------------------------------
% \chapter{Einleitung} % (fold)
% \label{cha:einleitung}

% chapter einleitung (end)
%--------------------------------------------------------------------------------

\chapter{Theoretische Grundlagen} % (fold)
\label{cha:theoretische_grundlagen}

\section{Beispiele} % (fold)
\label{sec:beispiele}

% section beispiele (end)

\subsection*{Ein Diagramm} % (fold)
\label{sub:ein_diagramm}

\begin{tikzpicture}
	\begin{axis}[]
		\addplot table[x index=0, y index=1]{Data/tabelle.txt};
	\end{axis}
\end{tikzpicture}
% subsection ein_diagramm (end)
   
\subsection*{Chemische Formeln} % (fold)
\label{sub:chemische_formeln}

    \ce{H2O} \\
    \ce{H2_{(aq)}}\\
    \ce{$A$ ->[\ce{+H2O}] $B$}\\
    \chemfig{H_3C-{{(CH_2)}_2}-CH_3}\\
    \chemfig{H_3C-C(=[1]O)-[7]O^\ominus}\\
    \chemname{\chemfig{*6(-(<[::-120]H)(*6(-(-[::-20]H_3C)%
(-[::-70]H_3C)-O-(*6(-=(--[:30]-[:-30]-[:30]-%
[:-30]CH_3)-=(-OH)-=))--(<:[::-120]H)-))--=(-CH_3)--)}}{THC}\\
% subsection chemische_formeln (end)


\subsection*{Eine Tabelle} % (fold)
\label{sub:eine_tabelle}
\begin{tabular}{l|rrr}
	 & Spalte 1 & Spalte 2 & Spalte 3 \\
	\hline
	erste Zeile & 1 & 2 & 3 \\
	zweite Zeile & 1 & 2 & 3 \\

\end{tabular}
% subsection eine_tabelle (end)

\subsection*{Zitate} % (fold)
\label{sub:zitate}


Ich zitiere hier \cite{gaukel2004untersuchungen} und \cite{tscheuschner2004grundzuge}. Die Quellen findet man am Ende der Arbeit im Literaturverzeichnis. Erstellt mit BibDesk.
% subsection zitate (end)


% chapter theoretische_grundlagen (end)
%--------------------------------------------------------------------------------

\chapter{Material und Methoden} % (fold)
\label{cha:material_und_methoden}

% chapter material_und_methoden (end)
%--------------------------------------------------------------------------------

\chapter{Ergebnisse und Diskussion} % (fold)
\label{cha:ergebnisse_und_diskussion}

% chapter ergebnisse_und_diskussion (end)
%--------------------------------------------------------------------------------

\chapter{Zusammenfassung und Ausblick} % (fold)
\label{cha:zusammenfassung_und_ausblick}

% chapter zusammen_und_ausblick (end)

%------------------------------------------------------------------------------
%	LITERATURVERZEICHNIS - ANHANG
%------------------------------------------------------------------------------
\bibliographystyle{../chicago-de}
\bibliography{../literatur}

\appendix

\listoffigures

\listoftables
%--------------------------------------------------------------------------------
\end{document}
